\section{OpenAI売上の500億ドル・700億ドル報道差}
OpenAIの9月末の年換算売上が約500億ドルだったとする報道と、従来伝えられた約700億ドルとの差が議論された。投稿者はFT記事を引用し、別の投稿では総額と純額の計上差だと説明している。Grok収集ではFT本文やOpenAI公式説明を確認できず、Reutersの匿名関係者に基づく報道も参照情報にとどまる。株価下落との因果や両数字の定義は未確定である。
\dailyxsource{Steve Hanke (@steve\_hanke)}{https://x.com/steve_hanke/status/2108629820229751113}

\section{AI生成数学原稿をめぐる批判が拡散}
集約アカウントが、OpenAIのAI生成数学原稿722本について、NYUのTristan Buckmaster氏が若手数学者のキャリアを損なうと批判した、と投稿した。原稿公開は観測窓外とされ、窓内では発言の引用と研究職への影響をめぐる反応が広がった。ただし引用元とされたCNBC映像は未視聴で、発言の正確さ、原稿の撤回件数、外部検証結果はいずれも確認されていない。
\dailyxsource{Big Brain AI (@realBigBrainAI)}{https://x.com/realBigBrainAI/status/2108513103734530130}

\section{Claude利用規約の改定をめぐる反応}
Anthropicが11月12日からClaudeへの「持続的で不必要な」残酷な扱いを禁じ、会話終了の対象にするとする投稿が続いた。通常の苛立ちや創作は対象外とする説明も流通している。今回は窓内の公式投稿と規約本文を確認できておらず、モデル福祉に関する議論や他社幹部の反応も二次投稿に依存する。規約の正確な適用範囲は未確認として扱う。
\dailyxsource{Dialed (@Notorious\_corbz)}{https://x.com/Notorious_corbz/status/2108676427507789955}

\section{Nace AI、Drex 1.5の公開を告知}
Nace AIは、意思決定向け小型モデルDrex 1.5をオープンソース化したと投稿した。128kコンテキスト、ローカル・クラウド利用、OllamaやOpenCode等との互換を訴求している。Decision Indexで他モデルを上回るという数値も引用投稿にあるが、独立再現はない。公開先の短縮URLは未展開で、ウェイトの実在、ライセンス、パラメータ数、評価条件は確認できていない。
\dailyxsource{Nace AI (@NaceAI)}{https://x.com/NaceAI/status/2108668596276867118}

\section{Nace AI、文書処理モデルNDI 1.0を紹介}
同社はコーディングエージェント向けの文書処理モデルNDI 1.0も紹介した。MCP・SDK・CLI対応、50言語、1,500万件超の金融ファイルによる学習を掲げ、精度やコストで既存モデルより優れると主張する。開発者向けページのURLは示されたものの本文は未取得で、Parse Index順位、GPT-Astra級精度、90％安価といった比較値の評価条件は確認されていない。
\dailyxsource{Nace AI (@NaceAI)}{https://x.com/NaceAI/status/2108420944146629080}

\section{ElevenLabs、合成音声検出を告知}
ElevenLabsは、企業への通話にAIエージェントが増えていることを背景に、ElevenAgents向けの合成音声検出機能を発表した。あわせてPersonal Agent Protocolの作業部会に参加するとしている。人間・企業・悪意ある主体が運用するエージェントを区別する必要性が示唆されるが、検出アルゴリズム、誤検知率、作業部会の仕様や参加者は投稿からは分からない。
\dailyxsource{ElevenLabs (@ElevenLabs)}{https://x.com/ElevenLabs/status/2108666775441461579}

\section{PA Score、言い換えと不確実性を評価}
Jiayi Xin氏は、言い換えによってLLMの不確実性が不当に変化する問題に対し、conformal scoreの設計で取り組むPA Scoreを紹介した。NeurIPS 2026採択を著者が報告し、arXiv、コード、ブログの所在も示している。確率的な信頼度を入力表現の揺れから切り離す試みとして注目される一方、論文要旨・実験設定や採択トラックの確認は今回の収集範囲外である。
\dailyxsource{Jiayi Xin (@RainaXin)}{https://x.com/RainaXin/status/2108584762621472840}

\section{Inference Auctions、推論資源に経済設計}
BerkeleyのNika Haghtalab氏は、LLM推論プラットフォームにオークションの仕組みを取り入れるInference Auctionsを紹介した。共著者の窓内投稿では、利用者の予算制約を考慮して入札を動的に調整するautobidderを提案したと説明する。推論サービングの計算資源配分と経済的誘因を一体で考える研究だが、論文URL、方式の詳細、既存方式との比較は未確認である。
\dailyxsource{Nika Haghtalab (@nhaghtal)}{https://x.com/nhaghtal/status/2108652254161215845}

\section{PLUTO、自然言語から航空軌道を最適化}
StanfordとNVIDIAの研究者Marco Pavone氏は、自然言語の要求を制約条件に変換し、逐次凸計画によって航空軌道を最適化するPLUTOを紹介した。AIが問題の定式化を担い、数値最適化手法が解を求める役割分担を強調している。研究者本人による窓内投稿だが、論文の短縮URLは未展開で、数値バックボーンの詳細や評価結果は未確認。コードも公開予定段階とされる。
\dailyxsource{Marco Pavone (@drmapavone)}{https://x.com/drmapavone/status/2108630478559551686}

\section{Maritime、エージェント専用の実行環境}
Maritimeの共同創業者は、エージェントごとに専用のコンピュータを割り当て、スタートアップが展開・管理できるクラウド基盤を紹介した。ホスティング料金は1エージェント月1ドルからと投稿している。エージェントを多数運用する際の分離と管理を訴求するが、数十億規模への拡張性、環境の隔離方式、契約条件や対応ランタイムは確認されておらず、価格も投稿者の主張として扱う。
\dailyxsource{Maria Gorskikh (@mariagorskikh)}{https://x.com/mariagorskikh/status/2108632744838488072}

\section{Agent Arena、長時間セッションを評価}
Arenaは、単発のペアワイズ投票だけでは捉えにくいエージェントの実務能力を、長時間の実セッションに基づいて評価するAgent Arenaを紹介した。5つの信号を因果追跡で測るというが、信号の定義や評価対象モデルは今回未取得である。資金調達などを扱う周辺ダイジェストは観測窓外の発表を混在させている可能性があり、本号は窓内の公式手法紹介に限定する。
\dailyxsource{Arena.ai (@arena)}{https://x.com/arena/status/2108674951624724970}

\section{Haiku 5.5、実タスクでコスト比較}
Claude Haiku 5.5の発表自体は観測窓外だが、窓内ではMinecraft用HTML生成や要約ジョブを使った費用比較が投稿された。個人検証では0.10ドル未満の生成例や、同じリスト価格でもトークン数・思考トークンの差からGPT-6 Lunaの方が安かったという主張がある。公式価格表や再現プロンプトは未取得で、ハーネスの違いも切り分けられていないため、一般的な価格優位を示す結果とは言えない。
\dailyxsource{Luckey Faraday (@luckeyfaraday)}{https://x.com/luckeyfaraday/status/2108678666477785278}

\section{企業向けGemini agentの解説が続く}
窓内の解説投稿は、Googleの企業向けGemini agentについて、Workspaceアカウント、監査主体、タスク別のモデル切替などを論じた。発表本体は10月8日頃とされ観測窓外の可能性が高く、今回はGoogle公式の当該投稿も確認できていない。GmailやDriveとの連携、支出上限、他社モデルへのルーティングなどは二次的な整理にとどまり、実際の提供範囲と権限管理仕様は未確認である。
\dailyxsource{HuangMing (@Zhm20220917)}{https://x.com/Zhm20220917/status/2108678880773234781}

\section{Saluki 27B、量子化後の性能が話題}
Qwen 3.8 27Bを基にしたUnderdog Salukiについて、54GBから7.89GBへの小型化、ツール呼び出し性能、Apache 2.0などを窓内の解説投稿が再掲した。ただし元の企業発表は10月8日で観測窓外。保持率96％や元モデルを上回るという比較は企業側の主張で、ウェイトURL、量子化方式、評価条件は未確認である。数学・推論性能の低下を懸念する投稿もある。
\dailyxsource{Md Ismail Sojal (@0x0SojalSec)}{https://x.com/0x0SojalSec/status/2108675821343027212}

\section{XRPLのエージェント決済、1300万件との主張}
RippleXのエンジニアリング責任者と名乗る投稿者は、XRPL上のエージェント取引が1,300万件を超え、x402加盟店166、台帳活動の約25％がエージェント決済だと述べた。エージェント間の支払いが実運用へ広がる可能性を示す数字だが、計測ダッシュボードや「エージェント取引」の定義は提示されていない。実績値ではなく投稿者の自己申告として記録する。
\dailyxsource{J. Ayo Akinyele (@ja\_akinyele)}{https://x.com/ja_akinyele/status/2108670853831671918}

\section{OpenAI・Anthropicの将来売上予測が流通}
Fireside Alphaは、SemiAnalysisのDylan Patel氏がOpenAIとAnthropicの2027年末までの売上について大きな規模を見込んだとする動画を投稿した。6,000億～7,000億ドル、場合によっては1.5兆ドルという数字が引用されるが、発言の全文脈や内訳は確認されていない。これは企業開示ではなく将来予測であり、足元の年換算売上を扱うトピック1とは対象期間も異なる。
\dailyxsource{Fireside Alpha (@firesidealpha)}{https://x.com/firesidealpha/status/2108649430534836413}
